\section{金属の強化機構}

\subsection{令和2年度 問8}
\label{q:r02-reg-applied-08}

\noindent\textbf{原典：} sources/R02/17\_applied\_science.pdf、PDF物理3頁。

\textbf{公式正答：} 5（single）。

\paragraph{問題の要旨} 金属の強化機構を識別し、焼入れに関する変態名の誤りを見抜く。

\paragraph{問われる核心} 材料物性

\paragraph{必要知識} 転位、加工硬化、固溶・析出・粒界強化、マルテンサイト変態

\paragraph{考え方（査読待ち初稿）} 鋼の焼入れ硬化で中心となるのはマルテンサイト変態である。⑤の「バウシンガー変態」は原典の誤答として保存し、解説中で正しい用語と区別する。したがって公式正答は⑤。

\paragraph{他の選択肢（査読待ち初稿）} ①は塑性加工による硬化、②は固溶による強化、③は温度依存の固溶限を利用する析出強化、④は結晶粒微細化の効果を述べる。転位運動を妨げる強化と、相変態を使う焼入れを対応付けて整理する。

\par 参考：\cite{martensite_cambridge}。

\par 難易度 3（校準初稿）、優先度未評価。
